
\section{\trHyperelliptic\trCurve}

接下来，我们讨论一类重要的\trCurve，即\trHyperelliptic\trCurve。在本节中，\emph{我们假设 $k$ 是\trCharacteristic{}不为 $2$ 的\trAlgClosedField}。（这些假设都可以放宽，但要付出一些代价。）

一条（射影、正则、整的）\trGenus{} $g$ \trCurve{} $C$ 称为\textbf{\trHyperelliptic}的（hyperelliptic），若它容许一个 $\mathbb{P}^1_k$ 的\trDoubleCover（即一个到 $\mathbb{P}^1_k$ 的、次数为 $2$ 的、必然有限的\trMorphism）。（关于从\trCurve{}到\trCurve{}的映射的次数的讨论，见 \S16.3.4。）为方便起见，当我们说 $C$ 是\trHyperelliptic{}的时，我们心目中便隐含着一个\trDoubleCover $\pi\colon C\to\mathbb{P}^1$ 的\emph{选取}。（我们日后会看到，若 $g\ge 2$，则这样的\trDoubleCover{}至多只有一个，见\mbox{命题~19.5.7}，所以这并不是一个多大的假设。）映射 $\pi$ 称为\textbf{\trHyperelliptic{}映射}（hyperelliptic map）。

由\mbox{练习~16.3.D}，$\mathbb{P}^1$ 的任一\trClosedPoint{} $p$ 的原像要么由一个点组成，要么由两个点组成。若 $\pi^{-1}(p)$ 是单个点，我们就说 $p$ 是一个\textbf{\trBranchPoint}（branch point），而 $\pi^{-1}(p)$ 是 $\pi$ 的一个\textbf{\trRamification{}点}（ramification point）。（\trRamification{}这一概念将在 \S21.7 中以更一般的方式定义。）

\textbf{19.5.1.\ \trhTheorem（\trHyperelliptic{}\trRiemann--\trHurwitz{}公式）。}\ ——\ \emph{设 $k=\bar{k}$ 且 $\operatorname{char}k\ne 2$，而 $\pi\colon C\to\mathbb{P}^1_k$ 是由 $k$ 上一条射影、正则、不可约、\trGenus{}为 $g$ 的\trCurve{}给出的\trDoubleCover。则 $\pi$ 有 $2g+2$ 个\trBranchPoint。}

这是\trRiemann--\trHurwitz{}公式的一个特例；我们将在 \S21.4 中陈述并证明该公式。你或许早已听说过：\trGenus{} $1$ 的复曲线会\trDoubleCover $\mathbb{P}^1$，并且在四个点上\trRamification。

为了证明\mbox{定理~19.5.1}，我们先证明下面这条。

\textbf{19.5.2.\ \trhProposition。}\ ——\ \emph{设 $\operatorname{char}k\ne 2$ 且 $k=\bar{k}$。给定 $\mathbb{P}^1$ 中 $r$ 个互不相同的点 $p_1,\dots,p_r$，则当 $r$ 为偶数时，恰有一个\trDoubleCover{}恰在这些点处\trRamification；而当 $r$ 为奇数时，则一个也没有。}

\emph{证明。}\quad 在 $\mathbb{P}^1$ 中取两点 $0$ 与 $\infty$，它们都不同于这 $r$ 个\trBranchPoint。全部\trBranchPoint{}都落在 $\mathbb{P}^1-\infty=\mathbb{A}^1=\Spec k[x]$ 之中。假设我们有 $\mathbb{A}^1$ 的一个\trDoubleCover $C'\to\mathbb{A}^1$，其中 $x$ 是 $\mathbb{A}^1$ 上的\trCoordinate。这便诱导出一个 $k(x)$ 上的二次\trFieldExtension{} $K$。由于 $\operatorname{char}k\ne 2$，这个扩张是\trGalois{}的。令 $\sigma\colon K\to K$ 为那个\trGalois{}\trInvolution。取 $K$ 中一个满足 $\sigma(y)=-y$ 的非零元素 $y$，于是 $1$ 与 $y$ 构成 $K$ 在域 $k(x)$ 上的一组基，并且二者都是 $\sigma$ 的特征向量。现在 $\sigma(y^2)=y^2$，故 $y^2\in k(x)$。我们可以用 $y$ 的一个适当的 $k(x)$ 倍替代 $y$，使得 $y^2$ 是一个无重因式的首一多项式。（我们正是在这里用到了 $k$ \trAlgClosed{}这一假设，为的是得到首项系数 $1$。）

于是 $y^2=x^N+a_{N-1}x^{N-1}+\cdots+a_0$，其中右边的多项式（记作 $f(x)$）无重根。正则性的\trJacobianCriterion（以\mbox{练习~13.2.F} 的面目出现）蕴涵：$\mathbb{A}^2=\Spec k[x,y]$ 中的这条\trCurve{} $C'_0$ 是\trRegular{}的。于是 $C'_0$ 是\trNormal{}的，并且与 $C'$ 具有相同的\trFunctionField。从而 $C'_0$ 与 $C'$ 都是 $\mathbb{A}^1$ 在由 $y$ 生成的那个域的有限扩张中的正规化，故二者\trIsomorphic。这样，我们便以一个显式的方程认出了 $C'$。

\trBranchPoint{}对应于那些使得 $y$ 恰有一个值的 $x$ 值，也就是说，对应于 $f(x)$ 的根。特别地，$N=r$，且
\[
f(x)=(x-p_1)\cdots(x-p_r),
\]
其中诸 $p_i$ 被理解为 $\bar{k}$ 中的元素。

在掌握了 $\mathbb{A}^1$ 上的情形之后，我们回到 $\mathbb{P}^1$ 上的情形。我们将考察这个\trBranchedCover{}在仿射开集 $\mathbb{P}^1\setminus\{0\}=\Spec k[u]$ 上的情形，其中 $u=1/x$。把上面的论证施加于 $\Spec k[u]$ 而非 $\Spec k[x]$，便表明任何这样的\trDoubleCover{}都必形如
\[
\begin{aligned}
C''&=\Spec k[Z,u]\big/\bigl(Z^2-(u-1/p_1)\cdots(u-1/p_r)\bigr)\\
&=\Spec k[Z,u]\Big/\Bigl(\bigl((-1)^r\textstyle\prod p_i\bigr)Z^2-u^rf(1/u)\Bigr)\\
&\longrightarrow \Spec k[u]=\mathbb{A}^1.
\end{aligned}
\]
因此，若在整个 $\mathbb{P}^1$ 上存在一个\trDoubleCover，则它必定是把 $C''$ 粘合到 $C'$ 上而得到的——这次粘合架在 $\Spec k[x]$ 与 $\Spec k[u]$ 粘合成 $\mathbb{P}^1$ 的那次粘合「之上」。

于是在 $K(C)$ 中，必有
\[
z^2=u^rf(1/u)=f(x)/x^r=y^2/x^r
\]
（其中 $z$ 是用 $(-1)^r\prod p_i$ 的一个平方根去乘 $Z$ 而得到的），由此 $z^2=y^2/x^r$。

若 $r$ 为偶数，把 $K(C)$ 看作由 $y$ 与 $x$ 生成的，则 $z$ 有两个可能的值：$z=\pm y/x^{r/2}$。必要时把 $z$ 改名为 $-z$ 之后，把这两块粘起来的方式便只有一种（我们取正的那个平方根）。

若 $r$ 为奇数，则结论由下面的\mbox{练习~19.5.A} 得出。

\textbf{19.5.A.（练习）}\quad 设 $\operatorname{char}k\ne 2$。证明：$x$ 在域 $k(x)[y]/(y^2-f(x))$ 中没有平方根，其中 $f$ 是一个根 $p_1,\dots,p_r$ 皆非零的多项式。（一个可能的提示：为什么 $\sqrt{3}\notin\mathbb{Q}(\sqrt{2})$？）请解释这如何证明 $r$ 为奇数时的\mbox{命题~19.5.2}。

\begin{flushright}
$\square$
\end{flushright}

为日后参考计，我们在此把\trHyperelliptic{}覆盖 $C\to\mathbb{P}^1$ 的这个显式的（双仿射）描述汇总如下。
\begin{equation*}\tag{19.5.2.1}\label{eq:hyperelliptic-two-affine}
\begin{tikzcd}[row sep=3.2em]
\Spec k[x,y]/(y^2-f(x)) \dar \rar["z=y/x^{r/2}","y=z/u^{r/2}"'] & \Spec k[u,z]/(z^2-u^rf(1/u)) \dar \\
\Spec k[x] \rar["u=1/x","x=1/u"'] & \Spec k[u]
\end{tikzcd}
\end{equation*}

\textbf{19.5.3.}\ \emph{若 $k$ 不\trAlgClosed。}\quad 若 $k$ 不\trAlgClosed{}（但\trCharacteristic{}不为 $2$），则上面的论证表明：若我们有 $\mathbb{A}^1$ 的一个\trDoubleCover，则它必形如 $y^2=af(x)$，其中 $f$ 是首一的（且 $a\ne 0$）。此外，若 $a$ 与 $a'$ 相差（在乘法意义下）一个 $(k^\times)^2$ 中的元素，则 $y^2=af(x)$ \trIsomorphic{}于 $y^2=a'f(x)$。你或许能利用这一点证明：（在 $k^\times\ne(k^\times)^2$ 的假设下）\trDoubleCover\emph{并不}由其\trBranchPoint{}决定。而且，这一失效由 $k^\times/(k^\times)^2$ 所分类。于是，我们有大量的\trCurve，它们在 $k$ 上互不\trIsomorphic，但在 $\bar{k}$ 上却变得\trIsomorphic。它们常被称作彼此的\emph{扭曲}（twists）。

（特别地，一旦我们定义了\trEllipticCurve，你将能够证明：$\mathbb{Q}$ 上存在两条具有相同的 $j$-不变量却互不\trIsomorphic{}的\trEllipticCurve；见\mbox{练习~19.10.H}。）

\textbf{19.5.4.}\ \emph{回到\trHyperelliptic{}\trRiemann--\trHurwitz{}公式（\mbox{定理~19.5.1}）的证明。}\quad 我们对 $\mathbb{P}^1$ 的那个在 $r$ 个不同点处\trRamification{}的唯一\trDoubleCover{}的显式描述，将使我们得以算出\trGenus，从而完成\mbox{定理~19.5.1} 的证明。

我们沿用\mbox{命题~19.5.2} 之证明中的记号 \eqref{eq:hyperelliptic-two-affine}。设 $\mathbb{P}^1$ 照例由 $\Spec k[x]$ 与 $\Spec k[u]$ 组成仿射覆盖，其中 $u=1/x$。设 $C\to\mathbb{P}^1$ 是一个\trDoubleCover，在 $\Spec k[x]$ 上由 $y^2=f(x)$ 给出，其中 $f$ 的次数为 $r$，且 $z^2=u^rf(1/u)$。那么 $C$ 有一个由 $\Spec k[x,y]/(y^2-f(x))$ 与 $\Spec k[u,z]/(z^2-u^rf(1/u))$ 组成的仿射开覆盖。$\OO_C$ 相应的\trCech\trComplexN{}为
\[
\begin{aligned}
0\longrightarrow k[x,y]/(y^2-f(x))\times k[u,z]/(z^2-u^rf(1/u))\overset{d}{\longrightarrow}{}\\
(k[x,y]/(y^2-f(x)))_x\longrightarrow 0.
\end{aligned}
\]
复形的第二个（也是最后一个）非零部分有一组由单项式 $x^ny^\epsilon$ 构成的基，其中 $n\in\mathbb{Z}$，而 $\epsilon=0$ 或 $1$。为计算\trGenus{} $g=h^1(C,\OO_C)$，我们必须计算 $\coker d$。我们可以用第一个因子 $k[x,y]/(y^2-f(x))$ 打到单项式 $x^ny^\epsilon$，其中 $n\in\mathbb{Z}^{\ge 0}$，而 $\epsilon=0$ 或 $1$。第二个因子的像由形如 $u^mz^\epsilon$ 的元素生成，其中 $m\ge 0$，而 $\epsilon=0$ 或 $1$。但 $u^mz^\epsilon=x^{-m}(y/x^{r/2})^\epsilon$。检视即知，余核有一组由单项式 $x^{-1}y$、$x^{-2}y$、……、$x^{-r/2+1}y$ 生成的基，从而其维数为 $r/2-1$。于是 $g=r/2-1$，\mbox{定理~19.5.1} 随之得证。

\begin{flushright}
$\square$
\end{flushright}

\textbf{19.5.5.}\ \emph{各种\trGenus{}的\trCurve。}\quad 作为\trHyperelliptic{}\trRiemann--\trHurwitz{}公式（\mbox{定理~19.5.1}）的一个推论，我们看到：在\trCharacteristic{}不为 $2$ 的\trAlgClosedField{}上，每个\trGenus{} $g\ge 0$ 都有\trCurve：要得到一条\trGenus{}为 $g$ 的\trCurve，只需考虑在 $2g+2$ 个互不相同的点处\trRamification{}的那个\trBranchedCover。唯一的那条\trGenus{} $0$ \trCurve{}就属于这种形式；我们将在 \S19.6.2 中看到，每条\trGenus{} $2$ 的\trCurve{}也都属于这种形式。我们马上就会看到，每条\trGenus{} $1$ 的\trCurve{}（提醒：在\trAlgClosedField{}上！）也是如此（\S19.9.5）。但指望一切\trCurve{}都属于这种形式就想得太多了；我们马上就会看到（\S19.7.4），存在\emph{不}是\trHyperelliptic{}的\trGenus{} $3$ \trCurve，而且我们甚至会得到「大多数」\trGenus{} $3$ \trCurve{}都不是\trHyperelliptic{}的这番启发式证据。日后我们还会给出（可以使之精确化的）模糊证据，表明当 $g>2$ 时「大多数」\trGenus{} $g$ \trCurve{}都不是\trHyperelliptic{}的（\S19.8.4）。

\textbf{19.5.B.（练习）}\quad 验证：一条\trGenus{}至少为 $1$ 的\trCurve{} $C$ 容许一个到 $\mathbb{P}^1$ 的次数 $2$ 覆盖，当且仅当它容许一个满足 $h^0(C,\mathscr{L})=2$ 的次数 $2$ \trInvertibleSheaf{} $\mathscr{L}$。在做这件事的过程中，你可能还愿意顺带验证：若 $C$ 是一条\trGenus{}至少为 $1$ 的\trCurve，而 $C$ 有一个带有至少 $2$ 个（线性无关的）截面的次数 $2$ \trInvertibleSheaf{} $\mathscr{L}$，则 $\mathscr{L}$ 恰有两个截面，并且这个 $\mathscr{L}$ 是\trBasePointFree{}的，从而给出一个\trHyperelliptic{}映射。

\textbf{19.5.6.\ \trhProposition。}\ ——\ \emph{设 $g\ge 2$。若 $\mathscr{L}$ 对应于一个\trHyperelliptic{}覆盖 $C\to\mathbb{P}^1$，则 $\mathscr{L}^{\otimes(g-1)}\cong\omega_C$。}

\emph{证明。}\quad 把\trHyperelliptic{}映射与第 $g-1$ 个\trVeroneseEmbedding{}复合：
\begin{equation*}\tag{19.5.6.1}\label{eq:hyperelliptic-veronese}
\begin{tikzcd}
C \rar["|\mathscr{L}|"] & \mathbb{P}^1 \rar["\nu_{g-1}\cong|\OO_{\mathbb{P}^1}(g-1)|"] & \mathbb{P}^{g-1}.
\end{tikzcd}
\end{equation*}
该复合对应于 $\mathscr{L}^{\otimes(g-1)}$。这个\trInvertibleSheaf{}的\trDegree{}为 $2g-2$（由\mbox{练习~18.4.G}）。

接下来我们证明拉回映射 $\mathrm{H}^0(\mathbb{P}^{g-1},\OO(1))\to\mathrm{H}^0(C,\mathscr{L}^{\otimes(g-1)})$ 是单射；也就是说，任一非零的 $s\in\mathrm{H}^0(\mathbb{P}^{g-1},\OO(1))$ 都拉回成 $\mathrm{H}^0(C,\mathscr{L}^{\otimes(g-1)})$ 中的一个非零元素；换言之，$C$ 在 $\mathbb{P}^{g-1}$ 中的像是非退化的。这是因为 $C$ 在 $\mathbb{P}^{g-1}$ 中的像正是 $\mathbb{P}^1$ 经 $\nu_{g-1}$ 的像（一条\trRationalNormalCurve；见 \S9.3.6）；而 $\mathbb{P}^1$ 在完全线性系 $|\OO_{\mathbb{P}^1}(g-1)|$ 之下的像是非退化的（\mbox{练习~15.2.C}）。

于是 $\mathscr{L}^{\otimes(g-1)}$ 至少有 $g$ 个截面。但由\mbox{练习~19.2.A}，\trDegree{}为 $2g-2$ 且带有（至少）$g$ 个截面的\trInvertibleSheaf{}只有\trCanonical{}层。

\begin{flushright}
$\square$
\end{flushright}

作为额外附赠，我们看到 \eqref{eq:hyperelliptic-veronese} 的复合正是\emph{完全}线性系 $|\mathscr{L}^{\otimes(g-1)}|$——$\mathscr{L}^{\otimes(g-1)}$ 的一切截面都由此而来。

\textbf{19.5.7.\ \trhProposition（\trGenus{} $\ge 2$ 的\trCurve{}至多以一种方式成为\trHyperelliptic）。}\ ——\ \emph{任何一条\trGenus{}至少为 $2$ 的\trCurve{} $C$ 都至多容许 $\mathbb{P}^1$ 的一个\trDoubleCover。更精确地说，若 $\mathscr{L}$ 与 $\mathscr{M}$ 是给出映射 $C\to\mathbb{P}^1$ 的两个次数 $2$ 的\trLineBundle，则 $\mathscr{L}\cong\mathscr{M}$。}

\emph{证明。}\quad 若 $C$ 是\trHyperelliptic{}的，则我们可以通过考虑由 $\omega_C$ 给出的\trCanonical\trLinearSystem{}找回那个\trHyperelliptic{}映射（即\textbf{\trCanonical{}映射}（canonical map），我们在接下来的几节中还要反复用到它）：它是到一个\trDegree{}为 $g-1$ 的\trRationalNormalCurve{}的\trDoubleCover（由前一条\trhProposition），而那条有理正规曲线\trIsomorphic{}于 $\mathbb{P}^1$。这个\trDoubleCover{}就是那个\trHyperelliptic{}覆盖（这同样由前一条\trhProposition{}的证明可知）。于是，\trHyperelliptic{}覆盖 $C\to\mathbb{P}^1$ 由\trCanonical{}映射所决定（故至多只有一个）。这个映射 $C\to\mathbb{P}^1$ 必须同时由 $\mathscr{L}$ 与 $\mathscr{M}$ 诱导，从而 $\mathscr{L}\cong\mathscr{M}$（利用把到\trProjectiveSpace{}的映射与\trLineBundle{}联系起来的\mbox{定理~15.2.2}）。

\begin{flushright}
$\square$
\end{flushright}

\textbf{19.5.C.（练习）}\quad 设 $C$ 是一条\trGenus{}至少为 $2$ 的\trHyperelliptic\trCurve。证明：\trHyperelliptic{}\trInvolution{}与 $C$ 的任何其他\trAutomorphism{} $g$ 交换。（提示：当你用 $g$ 去共轭这个\trInvolution{}时，会发生什么？）

\textbf{19.5.8.}\ \emph{「\trHyperelliptic\trCurve{}的空间」。}\quad 多亏了\mbox{命题~19.5.7}，我们现在可以分类\trGenus{}至少为 $2$ 的\trHyperelliptic\trCurve{}了。\trGenus{} $g\ge 2$ 的\trHyperelliptic\trCurve{}恰好对应于 $\mathbb{P}^1$ 上的 $2g+2$ 个互不相同的点——模掉 $S_{2g+2}$，也模掉 $\mathbb{P}^1$ 的\trAutomorphism。于是，「\trHyperelliptic\trCurve{}的空间」的维数为
\[
2g+2-\dim\Aut\mathbb{P}^1=2g-1.
\]
这并不是一个良定义的陈述，因为我们尚未严格地定义「\trHyperelliptic\trCurve{}的空间」——它是一个\trModuliSpace{}的一个例子。眼下，且把这当作一条合情合理的陈述。这个空间不可约且既约，这一点同样合情——它是某个不可约且既约之物的像。
